% !TeX program = xelatex
\documentclass[12pt,a4paper]{article}
\usepackage[margin=2cm]{geometry}
\usepackage{fontspec}
\usepackage{amsmath,amssymb,enumitem}
\setmainfont{Times New Roman}
\setlength{\parindent}{0pt}\setlength{\parskip}{6pt}
\begin{document}
\begin{center}{\Large\bfseries Đề thi Toán vào lớp 10 Hà Nội năm 2015\\Phần Đại số}\end{center}
\textbf{Câu 33.} Cho hai biểu thức
\[
P=\frac{x+3}{\sqrt{x}-2},\qquad Q=\frac{\sqrt{x}-1}{\sqrt{x}+2}+\frac{5\sqrt{x}-2}{x-4},
\quad x>0,\ x\ne4.
\]
\begin{enumerate}[label=\alph*)]
\item Tính $P$ khi $x=9$.
\item Rút gọn $Q$.
\item Tìm $x$ để $P/Q$ đạt giá trị nhỏ nhất.
\end{enumerate}
\textbf{Câu 34.} Một tàu tuần tra chạy ngược dòng $60$ km, sau đó chạy xuôi dòng $48$ km trên cùng một dòng sông có vận tốc dòng nước là $2$ km/giờ. Tính vận tốc của tàu khi nước yên lặng, biết thời gian xuôi dòng ít hơn thời gian ngược dòng $1$ giờ.

\textbf{Câu 35.} Giải hệ phương trình
\[
\begin{cases}
2(x+y)+\sqrt{x+1}=4,\\
(x+y)-3\sqrt{x+1}=-5.
\end{cases}
\]
\textbf{Câu 36.} Cho phương trình $x^2-(m+5)x+3m+6=0$ ($x$ là ẩn số).
\begin{enumerate}[label=\alph*)]
\item Chứng minh phương trình luôn có nghiệm thực với mọi $m\in\mathbb R$.
\item Tìm $m$ để hai nghiệm của phương trình là độ dài hai cạnh góc vuông của một tam giác vuông có cạnh huyền bằng $5$.
\end{enumerate}
\textbf{Câu 38.} Với $a,b\ge0$ thỏa mãn $a^2+b^2=4$, tìm giá trị lớn nhất của
\[
M=\frac{ab}{a+b+2}.
\]
\end{document}
